\chapter{Limitations and Future Work}

\section{Current Limitations}
\subsection{Dataset and Labeling}
\begin{itemize}[leftmargin=2em]
\item Class imbalance remains visible (for example, tiger vs boar support counts).
\item Background and lighting variation may still be limited compared to production use.
\item Label naming mismatch (hare vs rabbit aliasing) adds integration complexity.
\end{itemize}

\subsection{Model and Inference}
\begin{itemize}[leftmargin=2em]
\item RandomForest operates on per-frame static features and does not model temporal dynamics directly.
\item Confidence thresholds require manual tuning per runtime environment.
\item Similar hand poses can still produce occasional confusion.
\end{itemize}

\subsection{iOS Runtime and UX}
\begin{itemize}[leftmargin=2em]
\item Performance and stability may vary by device generation and camera quality.
\item Effects and tracking quality are sensitive to fast motion blur and partial occlusion.
\item Guidance timing and hold duration messaging should remain synchronized with logic constants.
\end{itemize}

\section{Addressed Since v1}
\begin{itemize}[leftmargin=2em]
\item Leak-free evaluation protocol honoring the Roboflow train/valid/test split (\texttt{\detokenize{train_unified_v2.py}}), with per-class reports and a documented augmentation study.
\item Unified-only CoreML model loading with loud failure instead of silent fallback to incompatible legacy models.
\item Rolling inference-latency instrumentation in \texttt{GestureRecognizer}.
\item Clock injection in \texttt{JutsuManager.processCandidate} and a 14-test SwiftPM suite (\texttt{swift\_tests/}) covering hold-to-commit, flicker rejection, label aliasing, sequence triggers, time limits, and battle-mode counters.
\end{itemize}

\section{Known Issues}
\begin{itemize}[leftmargin=2em]
\item \textbf{Kuchiyose is unreachable in free mode}: its sequence (boar, horse, monkey, bird) contains wind's (horse, monkey) as a substring; wind triggers at the monkey step and every trigger clears the accepted-sign history. Battle mode is unaffected because the counter follows the target sequence directly. Documented by \texttt{testKnownBugWindPreemptsKuchiyoseInFreeMode}; a fix requires a deliberate trigger-priority design decision (e.g., deferring shorter sequences while a longer one is in progress).
\end{itemize}

\section{Recommended Improvements}
\begin{enumerate}[leftmargin=2em]
\item Add temporal modeling (sequence-aware classifier or lightweight recurrent model).
\item Expand training data with multi-user and multi-environment captures.
\item Automate threshold calibration and per-class confidence analysis.
\item Add model versioning and reproducible experiment logging.
\item Build an automated evaluation suite for on-device regression testing.
\item Integrate screenshot and telemetry capture for documentation updates.
\end{enumerate}

\section{Documentation Maintenance Plan}
\begin{itemize}[leftmargin=2em]
\item Update metrics tables after each model retrain.
\item Replace placeholders with latest app screenshots before submission.
\item Keep jutsu-sequence table aligned with \texttt{AppModels.swift} source changes.
\end{itemize}
